\newpage
\section{VAN Bus}\label{sec:van_bus}

VAN (Vehicle Area Network) is a serial communication bus developed by the PSA group (Peugeot, Citroën) and Renault as an altenative to
the CAN bus. It's defined in the ISO standard \textbf{ISO 11519-3}. It utilizes the same physical layer as the CAN bus (a differential pair) and identical
bus arbitration method, but it differs in the way the data is packaged. Almost all of the information here was sourced from the Atmel TSS463C datasheet
\cite{tss463c_datasheet}.

\bigskip
\noindent
Like CAN, the VAN protocol uses identifiers to differentiate between messages and perform bus arbitration. Additionally, it also
has the ability for so called "In-frame" responses, bigger message sizes (28 bytes instead of 8 bytes in CAN) and the
ability to directly communicate with another device.

\bigskip
\noindent
Data on the VAN network is packed into frames. Frames consist of multiple parts: Start of frame (\textbf{SOF}),
message identifier (\textbf{IDEN}), command (\textbf{COM}), message data (\textbf{DATA}), frame checksum (\textbf{FCS}),
end of data signal (\textbf{EOD}), acknowledge signal (\textbf{ACK}) and the end of frame signal (\textbf{EOF}).

\begin{table}[H]
    \centering
    \fontsize{10}{10}
    \selectfont
    \begin{tabular}{|c|c|c|c|c|c|c|c|c|c|c|}
        \hline
        \multirow{2}{*}{\makecell{SOF             \\ 10 TS}}  & \multirow{2}{*}{\makecell{IDEN \\ 12 bit \\ 15 TS}} & \multicolumn{4}{c|}{\makecell{Command \\ (4 bit, 5 TS)}} & \multirow{2}{*}{\makecell{DATA max. \\ 28 byte \\ 280 TS}} & \multirow{2}{*}{\makecell{FCS \\ 15 bit \\ 18 TS}} & \multirow{2}{*}{\makecell{EOD \\ 2 TS}} & \multirow{2}{*}{\makecell{ACK \\ 2 TS}} & \multirow{2}{*}{\makecell{EOF \\ 4 TS}} \\[7mm] \cline{3-6}
         &  & EXT & RAK & R/W & RTR &  &  &  &  & \\
        \hline
    \end{tabular}
    \caption{VAN frame format}
    \label{fig:van_frame_format}
\end{table}

\subsection{Enhanced Manchester code}

All VAN frames are coded using Enhanced Manchester coding (e-Manchester).
With e-Manchester, after sending 3 NRZ bits, the fourth bit is sent as a Manchester bit.
Sending data this way ensures easier syncronization on the bus. This is a more predictable
method of bitstuffing as opposed to the CAN bus method.

In Table \ref{fig:van_frame_format} we can see the number of bits for each part of the frame, as well as
the number of Time Slots (TS). Since a Manchester bit takes two time slots to send, 4 bits in
a row encoded with e-Manchester takes 3 + 2 Time slots to send.

\begin{figure}[H]
    \centering
    \includegraphics[width=0.75\textwidth]{images/mux1-eman.pdf}
    \caption[]{e-Mancherster encoding \cite{mux1}}
    \label{fig:manchester_bits_mux1}
\end{figure}

\begin{figure}
    \centering
    \includegraphics[width=0.5\textwidth]{images/manchester_bits.pdf}
    \caption{NRZ and Manchester bits. \cite{tss463c_datasheet}}
    \label{fig:manchester_bits_tss463_docs}
\end{figure}

\subsection{Messages}

A device/module connected to the network can be one of three types:
\begin{itemize}
  \item The Autonomous module. A bus master. It can send \textbf{SOF} sequences, initiate data
        transfers and receive messages.
  \item The Synchronous access module. It cannot send \textbf{SOF} sequences, but can
        initiate data transfers and receive messages.
  \item The Slave module. It can only send messages using in-frame responses and can also
        receive messages.
\end{itemize}

\begin{figure}
  \begin{center}
    \includegraphics[width=0.75\textwidth]{images/ranking.pdf}
    \caption{Module types}
  \end{center}
\end{figure}

\noindent
The VAN protocol supports multiple message types. They are differentiated using the identifier (\textbf{IDEN})
and command (\textbf{COM}) fields. The identifier is 12 bits long and is followed by the 4 command bits. Due to the
way the bus arbitration works, a smaller identifier has priority on the bus. While the identifier value is completely
up to the user, the command bits are strictly defined and determine the message type that will be sent on the bus:
\begin{itemize}
    \item EXT - always 1
    \item RAK - „Request Acknowledge“. If 1, the receipient must acknowledge the message reception. If 0,
          the ACK signal must not be sent.
    \item R/W - „Read/Write“. Signals the data direction. If 0, then it's a `write' message, 
          a device is sending data intended for another device. If 1, then it's a `read' message,
          and it implies a request for data and expects an answer.
    \item RTR - „Remote Transmission Request“. If 0, the message contains data. If 1, the message
          contains no data and an in-frame response request is implied. The combination R/W = 0 and RTR = 1
          is forbidded on the bus.
\end{itemize}

\subsubsection{Message types}
\label{subsec:message-types}
Here we will list frame types relevant to this project. There are more, but we only use these.

\smallskip
\noindent
Reminder: Dominant = 0, Recessive = 1.

\paragraph{Normal frame}

The simplest frame type. One module sends an entire frame with data, and may or may not request an ACK.
In the case no ACK is requested, the message is treated as a broadcast message.

\begin{figure}[H]
    \centering
    \includegraphics[width=\textwidth]{images/normalframe_ack.pdf}
    \caption{Normal frame transimssion with ACK}
    \label{fig:normal_frame_ack}
\end{figure}

\begin{figure}[H]
    \centering
    \includegraphics[width=\textwidth]{images/normalframe_noack.pdf}
    \caption{Normal frame transimssion without ACK}
    \label{fig:normal_frame_noack}
\end{figure}

\paragraph{Reply Request Frame with Immediate Reply}
\label{p:immediate-reply}

This is the only frame in which a slave module can transmit data. The only difference on the
bus is the R/W bit changed state compared to the normal frame.
The bus master generates the \textbf{SOF}, the initiator generates the identifier and first 3 bits of the command,
and the ACK signal. The rest (RTR, data and checksum) are generated by the replying module.

\begin{figure}[H]
    \centering
    \includegraphics[width=\textwidth]{images/immediate_reply.pdf}
    \caption{Reply Request Frame with Immediate Reply}
    \label{fig:immediate_reply_frame}
\end{figure}

\subsection{VAN network}

We can connect multiple modules together to create a network. We can also connect multiple networks together
and use a gateway module to route data between networks. In the Peugeot 206 
we have 1 CAN network \emph{CAN Engine} and 3 VAN networks: \emph{VAN Comfort}, \emph{VAN Body 1} and \emph{VAN Body 2}.

\noindent
We can see the layout in figure \ref{fig:van_can_network}. From there we also see the gateway module, the BSI.
This module is present in all PSA group cars. This particular BSI is from the AEE2001 generation, the only generation
to use the VAN bus. Other generations, the AEE2004 and AEE2010 both exclusively use CAN. 
Our infotainment system will be connected in place of the CD changer, module number 8415 on figure \ref{fig:van_can_network}
on the \emph{VAN Comfort} bus.

\begin{figure}[H]
  \begin{center}
    \includegraphics[width=0.75\textwidth]{images/vannetwork.pdf}
    \caption{VAN and CAN Network layout in the Peugeot 206 \cite{sedre}}
    \label{fig:van_can_network}
  \end{center}
\end{figure}

